\begin{abstract}
Debiasing methods like Just Train Twice (JTT) implicitly assume that spurious features---shortcuts such as color or background cues---are learned earlier than causal features during gradient descent training, but this temporal ordering hypothesis has never been directly measured at the gradient level. We find that spurious features (color) converge 4 epochs earlier than core features (shape) during gradient descent---directly validating this assumption for the first time via ablation training and per-epoch gradient tracking on CMNIST (single-dataset proof-of-concept; generalization to Waterbirds, CelebA, NICO++ is future work). Temporal separation ($E_s=13$ vs $E_c=17$, $\Delta=4$ epochs) substantially exceeds the predicted 2-epoch threshold in proof-of-concept validation (full 10-seed statistical validation in progress). Early convolutional layers exhibit significantly higher spurious correlation than late layers ($\rho_j=0.003$ vs $0.001$, $p=0.0028$, Cohen's $d=0.25$), confirming that temporal gaps arise from architectural feature hierarchy where simpler low-level features stabilize before high-level semantics. Spurious-trained networks show 51\% lower forgetting rate (2.35 vs 4.82 events/sample), confirming stability advantages of simpler decision boundaries. However, gradient-aware training using CMNIST-derived $\rho_j$ values on Waterbirds failed catastrophically (39\% worst-group accuracy vs 86\% JTT target), revealing that neuron correlations are dataset-specific and cannot transfer across spurious feature types (color vs background). Our work provides the first mechanistic validation of temporal ordering underlying JTT/LfF reweighting methods while identifying cross-dataset $\rho_j$ transfer as a fundamental constraint for neuron-level interventions.
\end{abstract}
